本稿では, 6分布を対象とするモンテカルロ計算と, X 空間 / S 空間の区別に基づく理論整理を通じて, SLSC に基づく適合度評価を再検討した.
その結果, SLSC の標本サイズ依存性はより広い分布群でも確認され, 真の分布を当てはめた場合にも標本サイズの増加とともに低下した.
また, SLSC の値が変換に対して不変となるのは線形変換の場合に限られ, 非線形標準化では標準化関数に依存した距離評価となることを示した.
特に LN3 の S 空間 SLSC は, 真の分布が LN3 でない場合にも小さい値を与えやすく, モデル比較に影響しうる.

以上より, SLSC を用いる場合には, 固定閾値による一律判定に注意するとともに, 標準変量の定義が比較結果に与える影響を明示する必要がある.
分布間の適合度を共通尺度として比較したい場合には \(X\) 空間での評価を併用し, 非線形標準化を伴う SLSC を比較する場合には距離尺度の違いを踏まえて解釈することが望ましい.
今後は, LN3 そのものの統計的性質と, SLSC と併用可能な尤度・情報量規準に基づく分布選択法を検討する必要がある.
